\documentclass[11pt]{article}

\usepackage[a4paper,margin=1in]{geometry}
\usepackage{amsmath,amssymb,amsthm}

\newtheorem*{problem}{Problem}
\newtheorem*{theorem}{Theorem to prove}
\newtheorem*{remark}{Remark}

\DeclareMathOperator{\adj}{adj}
\DeclareMathOperator{\Tr}{Tr}

\newcommand{\Cat}{\operatorname{Cat}}

\begin{document}

\begin{center}
{\Large A self-contained problem for the Appendix C universality theorem}
\end{center}

\begin{problem}
This is a finite algebraic-combinatorial problem.  There are no random matrices
in the statement.  The goal is to prove the universal first-order threshold
theorem from the moment and Hankel data below, not merely to finish the last
divisibility step after the hard identities have already been supplied.

Work over the polynomial ring \(\mathbb Q[\lambda]\).  Let
\[
  \Cat(r)=\frac{1}{r+1}\binom{2r}{r}
\]
be the Catalan number.  Define the balanced moment polynomials
\[
  M_n(\lambda)
  =
  \sum_{r=0}^{\lfloor n/2\rfloor}
    \binom{n}{2r}\Cat(r)\lambda^{\,n-r},
  \qquad n\ge0.
\]
Equivalently, the generating series \(C(z)=\sum_{n\ge0}M_n(\lambda)z^n\)
is characterized by
\[
  C(z)-1=\lambda z C(z)\bigl(1+zC(z)\bigr).
\]

Define the first-defect moment polynomials \(M_n^{(1)}(\lambda)\) by
\[
\begin{aligned}
M_n^{(1)}
&=
\lambda
\sum_{\substack{a,b,c\ge0\\a+b+c+3=n}}
  (a+1)M_aM_bM_c  \\
&\quad+
\lambda
\sum_{\substack{a,b,c,d\ge0\\a+b+c+d+4=n}}
  (a+1)M_aM_bM_cM_d .
\end{aligned}
\]
The two sums are understood as \(0\) if their indexing sets are empty.
This definition is the rooted first-defect correction: exactly one exceptional
block has size \(3\) or \(4\), and all remaining pieces are balanced.

For \(m\ge0\), form the \((m+1)\times(m+1)\) Hankel matrices
\[
  H_m(\lambda)=\bigl(M_{i+j+1}(\lambda)\bigr)_{0\le i,j\le m},
  \qquad
  H_m^{(1)}(\lambda)=
    \bigl(M_{i+j+1}^{(1)}(\lambda)\bigr)_{0\le i,j\le m}.
\]
Set
\[
  f_{0,m}(\lambda)=\det H_m(\lambda),
  \qquad
  f_{1,m}(\lambda)=\Tr\!\bigl(\adj(H_m(\lambda))H_m^{(1)}(\lambda)\bigr).
\]
Thus \(f_{1,m}\) is the coefficient of \(\varepsilon\) in
\[
  \det\bigl(H_m(\lambda)+\varepsilon H_m^{(1)}(\lambda)\bigr).
\]

Finally define polynomials \(d_r(\lambda)\) by
\[
  d_0=1,\qquad d_1=\lambda,\qquad
  d_{r+2}=\lambda d_{r+1}-\lambda d_r\quad(r\ge0),
\]
and define \(g_m(\lambda)\) by removing the zero-root power from \(d_{m+1}\):
\[
  d_{m+1}(\lambda)=
  \lambda^{\lceil (m+1)/2\rceil}g_m(\lambda),
  \qquad g_m(0)\ne0.
\]

\begin{theorem}
For every \(m\ge0\), prove the following.

\[
  f_{0,m}(\lambda)
  =
  \lambda^{m(m+1)/2}d_{m+1}(\lambda).
  \tag{A}
\]

\[
  g_m(\lambda)\mid
  \bigl(f_{1,m}(\lambda)-f_{0,m}'(\lambda)\bigr).
  \tag{B}
\]

Consequently, if \(\alpha\ne0\) is a simple root of \(g_m\), then
\[
  f_{1,m}(\alpha)=f_{0,m}'(\alpha).
  \tag{C}
\]
Hence any scalar \(\kappa\) satisfying the first-order threshold equation
\[
  \kappa f_{0,m}'(\alpha)+f_{1,m}(\alpha)=0
  \tag{D}
\]
must be
\[
  \kappa=-1.
  \tag{E}
\]
\end{theorem}

The exercise is intentionally open-ended.  A solution may use determinants,
orthogonal polynomials, generating functions, non-crossing partitions,
continued fractions, linear algebra, resultants, explicit kernel vectors, or
any other purely algebraic method.  The choice of proof strategy is part of
the exercise.

The following shortcuts are not allowed as assumptions: the congruence
\((B)\), the equality \((C)\), a prepackaged trace identity for
\(f_{1,m}-f_{0,m}'\), or a prepackaged vanishing statement at the roots of
\(g_m\).  Such statements may of course be proved inside a solution.
\end{problem}

\begin{remark}
This is the pure-math core of Appendix C.  In the paper, the conclusion
\((B)\) says that the first finite-size correction has the same values as the
\(\lambda\)-derivative of the leading Hankel determinant at every leading
threshold.  Equation \((D)\) is the formal first-order root equation; once
\((C)\) is known, the universal coefficient \(-1\) follows immediately.
\end{remark}

\end{document}
